\section*{Executive Summary}
\addcontentsline{toc}{section}{Executive Summary}

On 1 January 2026 the European Union began charging importers for the carbon
embedded in imported steel goods. In the preceding reporting period Taiwan
shipped \textbf{3.74 million tonnes} of covered goods to the EU --- mostly
screws and fasteners produced by \textbf{about 2,600 small, family-run
factories} concentrated in the Gangshan district of Kaohsiung, the ``Kingdom of
Screws''\cite{tpt-cbam}. From this year each of those factories must either
deliver verified, product-level carbon data to its European buyer, or let that
buyer absorb punitive ``default values'' --- marked up 10\% now and 30\% by
2028\cite{omm} --- until the order moves to a supplier who can supply the data.

The information needed to avoid that penalty exists, but it is buried in a
2,400-page EU implementing regulation\cite{omm} that a senior factory owner,
whose back office is a family member with a spreadsheet, cannot practically
read --- and no tool in Taiwan's current public or commercial ecosystem
produces the required artifact for him\cite{caas,idb,tpt-cbam}. This is the
digital divide of the AI era in its most economically acute form: not the
absence of connectivity, but the inability to produce the data that now decides
who may participate in global markets.

\textbf{CarbonPass closes that divide from a photograph.} Inside LINE, in
Mandarin or Taiwanese, an owner photographs an electricity bill and a set of
invoices. A local-first AI --- raw data never leaves the factory --- returns
three outputs: (1) the verifier-ready, per-product carbon data pack the EU
buyer requires, in the European Commission's own template; (2) a decision
screen quantifying what the data is worth, \emph{``with your figures the buyer
pays $\approx$EUR~150 per tonne; without them, EUR~450--750 and
rising''}\cite{ecprice,specialinsert}; and (3) a production schedule, computed
from Taiwan's open 10-minute grid data\cite{taipowerdata}, that shifts flexible
loads to cheaper, lower-carbon hours --- reducing the electricity bill now and
the next data pack's emissions with it.

The proof of concept is built entirely on Taiwan's open data (MOENV emission
coefficients, Taipower grid telemetry and tariffs) and the EU's published
methodology, and it gives data back: a cleaned hourly grid-carbon-intensity
dataset and an anonymized sector benchmark, openly licensed. The beachhead
population is precisely defined --- 2,600 CBAM-exposed exporters within a
$\sim$1,800-firm, top-three world export cluster --- and the same engine scales
as the mechanism expands: the European Parliament votes in September 2026 on
extending coverage to $\sim$180 further product categories\cite{akin}, and,
beyond Taiwan, every SME-exporting economy faces the same wall.

\section{Team Information}\label{sec:team}

\emph{Application form field 1. Member names, nationalities, contact emails,
phone numbers, gender, and current country of residence are entered in the
online application form; the primary and secondary contacts are listed first,
as required.}

\begin{formfield}{Team composition and eligibility}
The team comprises between 3 and 10 members and includes at least one member
holding nationality other than that of the Republic of China (Taiwan), in
accordance with the eligibility rules\cite{handbook}. Two members are
designated as primary and secondary contacts. The team spans four
competencies matched to the work: machine learning and document intelligence;
energy-systems and optimization engineering; EU/Taiwan carbon-regulation
analysis; and fastener-industry domain knowledge. All project content,
documentation, and presentation materials are prepared and submitted in
English.
\end{formfield}

\begin{formfield}{Originality and commitments}
The entry is the original work of the team, created for this competition. It
has not been sold or awarded elsewhere, so the 50\% modification rule does not
apply. Intellectual-property arrangements among members follow the Handbook's
terms\cite{handbook}. The team commits to full participation in the mentorship
period, implementation-result presentations, and follow-up benefit tracking
required of advancing teams.
\end{formfield}
